% =====================================================================
%  Proposta — Campus Virtual Training & Campus LINGO
%  Stile ispirato a "For Dummies" — versione compatta 3 pagine
%  Compilare con: pdflatex proposta.tex
% =====================================================================
\documentclass[11pt,a4paper]{article}

% ---------- Pacchetti ----------
\usepackage[utf8]{inputenc}
\usepackage[T1]{fontenc}
\usepackage[italian]{babel}
\usepackage{microtype}
\usepackage[a4paper,top=1cm,bottom=1.5cm,left=2.5cm,right=2.5cm,marginparwidth=2cm,marginparsep=4mm,headheight=22pt]{geometry}
\usepackage{xcolor}
\usepackage{tikz}
\usepackage{tcolorbox}
\tcbuselibrary{skins,breakable}
\usepackage{titlesec}
\usepackage{fancyhdr}
\usepackage{enumitem}
\usepackage{lettrine}
\usepackage{marginnote}
\usepackage{hyperref}
\usepackage{parskip}

% ---------- Colori "For Dummies" ----------
\definecolor{dummiesyellow}{RGB}{255,204,0}
\definecolor{dummiesblack}{RGB}{0,0,0}
\definecolor{dummiesgray}{RGB}{90,90,90}
\definecolor{dummieslightyellow}{RGB}{255,243,191}

% ---------- Hyperref ----------
\hypersetup{
  colorlinks=true,
  linkcolor=dummiesblack,
  urlcolor=dummiesblack,
  citecolor=dummiesblack,
  pdftitle={Proposta Campus Virtual Training},
}

% ---------- Stile sezioni ----------
\titleformat{\section}
  {\normalfont\sffamily\large\bfseries\color{dummiesblack}}
  {\thesection}{1em}{}
  [{\color{dummiesyellow}\rule{\textwidth}{1.2pt}}]
\titlespacing*{\section}{0pt}{6pt}{2pt}

% ---------- Header e footer ----------
\pagestyle{fancy}
\fancyhf{}
\fancyfoot[C]{\sffamily\small\thepage}
\renewcommand{\headrulewidth}{0pt}

% ---------- Icone ----------
\newcommand{\iconTarget}{%
  \begin{tikzpicture}[scale=0.3]
    \fill[dummiesyellow] (0,0) circle (1);
    \draw[dummiesblack,thick] (0,0) circle (1);
    \fill[white] (0,0) circle (0.6);
    \draw[dummiesblack,thick] (0,0) circle (0.6);
    \fill[dummiesyellow] (0,0) circle (0.25);
    \draw[dummiesblack,very thick] (0.5,1.2) -- (0.05,0.1);
    \fill[dummiesblack] (0.5,1.2) -- (0.3,0.8) -- (0.7,0.9) -- cycle;
  \end{tikzpicture}}

\newcommand{\iconRemember}{%
  \begin{tikzpicture}[scale=0.3]
    \draw[dummiesblack,very thick] (-1.2,0) -- (1.2,0);
    \fill[dummiesyellow] (0,0) circle (0.45);
    \draw[dummiesblack,thick] (0,0) circle (0.45);
  \end{tikzpicture}}

\newcommand{\iconWarning}{%
  \begin{tikzpicture}[scale=0.3]
    \fill[dummiesyellow] (0,1.2) -- (-1.1,-0.7) -- (1.1,-0.7) -- cycle;
    \draw[dummiesblack,thick] (0,1.2) -- (-1.1,-0.7) -- (1.1,-0.7) -- cycle;
    \node[dummiesblack,font=\bfseries\small] at (0,0) {!};
  \end{tikzpicture}}

% ---------- Box ----------
\newcommand{\dummiestip}[2]{%
  \marginnote{\iconTarget\\\sffamily\footnotesize\bfseries #1}%
  \begin{tcolorbox}[colback=dummieslightyellow,colframe=dummiesyellow,
    boxrule=1pt,arc=2pt,left=6pt,right=6pt,top=3pt,bottom=3pt,breakable]
    \small #2
  \end{tcolorbox}}

\newcommand{\dummiesremember}[2]{%
  \marginnote{\iconRemember\\\sffamily\footnotesize\bfseries #1}%
  \begin{tcolorbox}[colback=dummieslightyellow,colframe=dummiesyellow,
    boxrule=1pt,arc=2pt,left=6pt,right=6pt,top=3pt,bottom=3pt,breakable]
    \small #2
  \end{tcolorbox}}

\newcommand{\dummieswarning}[2]{%
  \marginnote{\iconWarning\\\sffamily\footnotesize\bfseries #1}%
  \begin{tcolorbox}[colback=dummieslightyellow,colframe=dummiesyellow,
    boxrule=1pt,arc=2pt,left=6pt,right=6pt,top=3pt,bottom=3pt,breakable]
    \small #2
  \end{tcolorbox}}

% =====================================================================
\begin{document}
\thispagestyle{fancy}

% ---------- Intestazione compatta stile "For Dummies" ----------
\noindent
\begin{tikzpicture}
  \fill[dummiesyellow] (0,0) rectangle (\textwidth,2.2);
  \fill[dummiesblack] (0,-0.6) rectangle (\textwidth,0);
  \node[anchor=west,font=\sffamily\bfseries,text=dummiesblack]
    at (0.4,1.55) {\fontsize{22}{26}\selectfont Proposta};
  \node[anchor=west,font=\sffamily\bfseries,text=dummiesblack]
    at (0.4,0.75) {\fontsize{18}{22}\selectfont Un'evoluzione condivisa};
  \node[anchor=west,font=\sffamily\small,text=white]
    at (0.4,-0.3) {Campus Virtual Training \& Campus LINGO \textemdash{} Scm Campus};
\end{tikzpicture}

\vspace{0.4em}

\lettrine[lines=2,lhang=0.1,findent=2pt]{N}{egli ultimi anni}, lavorando
nella formazione, mi sono reso conto di quanto la tecnologia non sia
più solo un supporto, ma una leva concreta per far evolvere davvero
quello che facciamo.

E, un po' alla volta, ho iniziato a muovermi in questa direzione. Oggi
non mi sento più solo ``il formatore'': mi piace farlo, ma nel
frattempo sono nati \textbf{Campus Virtual Training} e
\textbf{Campus LINGO}. Sono strumenti già utilizzabili, con un
potenziale che secondo me va oltre l'uso interno.

\dummiesremember{Il punto}{%
Così come sono oggi, riesco a portarli avanti ``a pezzi'', nei
ritagli. Ma per trasformarli in qualcosa di \textbf{strutturato},
\textbf{scalabile} e anche \textbf{proponibile ai clienti}, serve un
passo in più.}
{\color{dummiesyellow}\rule{\textwidth}{1.2pt}}
\section*{Perché il tuo ruolo è fondamentale}


Io non sto cercando di cambiare direzione da solo, né di costruire
qualcosa ``a lato''. Al contrario, credo che questa possa essere
un'evoluzione naturale del lavoro che stiamo già facendo --- ma
guidata insieme, con una visione chiara e condivisa.

Mi piacerebbe costruire questo percorso con te:
\begin{itemize}[leftmargin=*,topsep=2pt,itemsep=1pt]
  \item dare una forma più concreta a questi progetti;
  \item capire come integrarli nel nostro perimetro;
  \item se ha senso, iniziare a strutturarli come veri e propri
        \textbf{prodotti}.
\end{itemize}

\dummiestip{In pratica}{%
Io porto l'iniziativa e lo sviluppo, ma ho bisogno di un contesto, di
una direzione e di un allineamento forte con te per farli crescere
davvero.}
{\color{dummiesyellow}\rule{\textwidth}{1.2pt}}
\section*{Un'estensione del lavoro che faccio}


La formazione resta il cuore di quello che faccio: è il terreno da cui
nascono queste idee e quello su cui continuano ad avere senso. Quello
che propongo è affiancarci un tassello in più, che amplia ciò che
portiamo ai clienti e l'impatto che possiamo avere come team.

\dummieswarning{Attenzione}{%
Questa non è una proposta di separazione o di cambio di rotta: è una
richiesta di allineamento per trasformare qualcosa che sta già
crescendo in un asset condiviso e sostenibile, integrato con l'attività
di formazione che continua.}
{\color{dummiesyellow}\rule{\textwidth}{1.2pt}}
\section*{Il passo successivo}


Se per te ha senso, mi piacerebbe ragionarci insieme e capire come
trasformare questa cosa in un percorso concreto.

\vspace{0.6em}
\begin{center}
\itshape Prometto che non ti porto solo idee\ldots{} ma anche un po' di
lavoro in più.
\end{center}

\end{document}
